
\section{Lecciones aprendidas para el PFC (SREI)}

El incidente de Cloudflare demuestra que una
configuración incorrecta, un despliegue sin controles
y la falta de reversión pueden afectar varios servicios
al mismo tiempo \cite{cloudflare2022}. En SREI
identificamos tres riesgos relacionados con las
tecnologías que utilizamos y las medidas que
incorporaremos en TA-PFC-E1.

\begin{table}[htbp]
\centering
\small
\begin{tabularx}{\linewidth}{@{}L L@{}}
\toprule
\textbf{Riesgo en SREI} & \textbf{Medida preventiva} \\
\midrule
Una configuración incorrecta de la conexión entre
Angular y Spring Boot podría impedir que los usuarios
accedan a las funciones del sistema.
&
Verificar las direcciones y configuraciones de conexión
en un entorno de pruebas antes de publicar cambios.
\\
\addlinespace
Una actualización del backend Spring Boot o de las
tablas de PostgreSQL podría provocar errores en el
registro y consulta de eventos.
&
Probar las actualizaciones y migraciones antes de
aplicarlas, conservando respaldos de la base de datos.
\\
\addlinespace
Un cambio en Docker Compose podría afectar
simultáneamente los contenedores y dejar SREI
fuera de servicio.
&
Probar los cambios antes del despliegue, actualizar
los servicios de forma controlada y conservar las
versiones anteriores para facilitar la reversión.
\\
\bottomrule
\end{tabularx}
\end{table}
